\documentclass[11pt]{article}
\usepackage[margin=1in]{geometry}
\usepackage{amsmath,amssymb,amsthm,mathtools}
\usepackage{microtype}
\usepackage{booktabs}
\usepackage{enumitem}
\usepackage[hidelinks]{hyperref}

\newtheorem{theorem}{Theorem}
\newtheorem{lemma}[theorem]{Lemma}
\newtheorem{proposition}[theorem]{Proposition}
\newtheorem{corollary}[theorem]{Corollary}
\newtheorem{remark}[theorem]{Remark}
\newtheorem{definition}[theorem]{Definition}

\newcommand{\DB}{\operatorname{DB}}
\newcommand{\normal}{\mathrm{normal}}
\newcommand{\Unif}{\operatorname{Unif}}
\newcommand{\1}{\mathbf{1}}

\title{\textbf{A Finite De Bruijn Certificate Improving the Upper Bound\\
for One-Round Randomized 2-Coloring of Cycles}}
\author{Ryutaro Yonezu\\Independent Researcher}
\date{September 9, 2026\\\small Preprint --- not peer reviewed}

\begin{document}
\maketitle

\begin{center}
\small
Paper DOI: \href{https://doi.org/10.5281/zenodo.22670656}{10.5281/zenodo.22670656}\\
Software DOI: \href{https://doi.org/10.5281/zenodo.22669456}{10.5281/zenodo.22669456}
\end{center}

\begin{abstract}
Flin, Raevskaya, Stimpert, Suomela, and Yang recently reduced one-round
randomized 2-coloring of directed cycles to cuts in three-dimensional De
Bruijn graphs and proved
\[
  0.23879\le p^*<0.24118.
\]
Their formally verified three-parameter construction has exact monochromatic-edge
probability
\[
  \frac{94835}{393216}=0.2411778767903646\ldots.
\]
We give an explicit finite coloring of the normal De Bruijn graph
$\DB_{\normal}(497)$ with exactly $14{,}714{,}994{,}317$ monochromatic edges
out of $61{,}013{,}446{,}081$ directed edges.  Hence
\[
  p^*\le
  \frac{14{,}714{,}994{,}317}{61{,}013{,}446{,}081}
  =0.2411762531404098\ldots
  <\frac{94835}{393216}.
\]
The coloring is the midpoint discretization of the published
three-parameter rule, followed by $4750$ explicitly listed vertex flips.
The strict improvement is certified by exact integer arithmetic; in
particular, the relevant cross-product difference is
$-38{,}953{,}738{,}163$.  We also show that no local flips are needed merely
to obtain a strict improvement: the unmodified midpoint discretization at
$n = 281$ already beats the published exact value.  The contribution is a
finite, machine-checkable upper-bound certificate within the De Bruijn
framework of Flin et al.; we do not claim a new general framework, a lower
bound, or optimality of the displayed finite cuts.
\end{abstract}

\noindent\textbf{Keywords.} distributed graph algorithms; randomized local algorithms; cycle coloring; De Bruijn graphs; maximum cut; finite certificates.

\section{Introduction}

Consider a directed cycle whose nodes run a randomized distributed algorithm
for one communication round and then output one of two colors.  A proper
2-coloring cannot always be produced, so the objective is to minimize the
expected fraction of monochromatic edges.  In the standard continuous-randomness
formulation, a one-round algorithm is a measurable function
\[
  f:[0,1]^3\to\{0,1\},
\]
where the three arguments are the independent random values held by a node's
predecessor, the node itself, and its successor.  If $A,B,C,D$ are i.i.d.
$\Unif[0,1]$, define
\[
  p(f):=\Pr\bigl[f(A,B,C)=f(B,C,D)\bigr],
  \qquad
  p^*:=\inf_f p(f).
\]

Flin et al.~\cite{flin2026arxiv,flin2026podc} developed a systematic finite
reformulation of this problem using normal and distinct three-dimensional De
Bruijn graphs.  Their main result is
\[
  0.23879\le p^*<0.24118.
\]
For the upper bound, they discovered a three-parameter recursive cutoff rule
and formally verified its exact value in Lean 4.  In the public formalization
at the pinned commit cited in the bibliography, the file
\texttt{Distributed2Coloring/UpperBound/Recursive3Param.lean} defines
\[
  t = \frac58,
  \qquad
  t_1 = \frac38,
  \qquad
  t_2 = \frac{17}{32},
\]
and the file
\texttt{Distributed2Coloring/UpperBound/Recursive3Param/Final.lean}
proves~\cite{flin2026lean}
\[
  p_0:=\frac{94835}{393216}
  =0.2411778767903646\ldots.
\]
The repository's \texttt{Distributed2Coloring/MainResults.lean} packages
the resulting public statement $p^*<0.24118$; it does not contain the exact
rational equality, which is proved in the \texttt{Recursive3Param} files
just cited.

This note asks a deliberately narrow question: can the same De Bruijn
framework provide a finite coloring whose exact ratio is strictly smaller
than $p_0$?  The answer is yes.

\paragraph{Main result.}
We construct an explicit coloring $h_{497}$ of $\DB_{\normal}(497)$ with
\[
  M_{497}=14{,}714{,}994{,}317
\]
monochromatic edges among
\[
  497^4=61{,}013{,}446{,}081
\]
directed edges.  By the De Bruijn upper-bound lemma of Flin et al., this gives
\[
  p^*\le\frac{M_{497}}{497^4}
  =0.2411762531404098\ldots.
\]
The comparison with $p_0$ is exact:
\[
  M_{497}\cdot393216
  -94835\cdot497^4
  =-38{,}953{,}738{,}163<0.
\]
Thus the new bound is not a floating-point artifact.

\paragraph{A simpler strict improvement.}
The local-refinement step is not essential for the qualitative statement.
The midpoint discretization of the published three-parameter rule on
$\DB_{\normal}(281)$, with no vertex flips, has exactly
$1{,}503{,}704{,}985$ monochromatic edges among $281^4=6{,}234{,}839{,}521$
edges, and
\[
  1{,}503{,}704{,}985\cdot393216
  -94835\cdot281^4
  =-146{,}592{,}275<0.
\]
Hence even this short deterministic finite discretization already yields a
strictly smaller certified upper bound than $p_0$.

\paragraph{Scope.}
The novelty claimed here is only the new finite upper-bound certificate and
its exact verification.  The reduction from one-round randomized cycle
coloring to De Bruijn cuts, the three-parameter base rule, and the idea of
searching for good finite cuts all come from Flin et al.~\cite{flin2026arxiv}.
We do not claim that $h_{497}$ is an optimal coloring of
$\DB_{\normal}(497)$, and we do not improve the lower bound $0.23879$.

\section{De Bruijn formulation}

For an integer $n\ge2$, write $[n]=\{0,1,\dots,n-1\}$.  The normal
three-dimensional De Bruijn graph is
\[
  \DB_{\normal}(n)=(V_n,E_n),
\]
where
\[
  V_n=[n]^3
\]
and
\[
  E_n=
  \left\{
    \bigl((a,b,c),(b,c,d)\bigr):a,b,c,d\in[n]
  \right\}.
\]
Therefore
\[
  |V_n|=n^3,
  \qquad
  |E_n|=n^4.
\]

For a coloring $h:V_n\to\{0,1\}$, let
\[
  M(h):=
  \#\left\{
    (a,b,c,d)\in[n]^4:
    h(a,b,c)=h(b,c,d)
  \right\}
\]
be its number of monochromatic directed edges, and put
\[
  \rho(h):=\frac{M(h)}{n^4}.
\]
Let $p_{\normal}(n)$ be the minimum of $\rho(h)$ over all such colorings.

We use the following consequence of Lemma~4.1 of Flin et al.~\cite{flin2026arxiv}.

\begin{lemma}[Finite De Bruijn certificate]
For every $n\ge2$ and every coloring $h:V_n\to\{0,1\}$,
\[
  p^*\le p_{\normal}(n)\le\rho(h).
\]
\end{lemma}

\begin{proof}
The first inequality is exactly the normal-De-Bruijn upper-bound lemma of
Flin et al.  The second follows from the definition of
$p_{\normal}(n)$ as the minimum monochromatic-edge fraction among all
colorings of $\DB_{\normal}(n)$.
\end{proof}

Thus a new upper bound on $p^*$ can be proved by exhibiting any explicit
finite coloring and counting its monochromatic edges exactly.

\section{Base coloring}

We start from the three-parameter recursive cutoff rule of Flin et
al.~\cite{flin2026lean}.  The parameter definitions below are a direct
transcription of
\texttt{Distributed2Coloring/UpperBound/Recursive3Param.lean} at the pinned
commit cited in the bibliography.
This section states the rule explicitly so
that the finite certificate does not depend on an informal geometric
description.
Set
\[
  t_1 = \frac38,
  \qquad
  t_2 = \frac{17}{32},
  \qquad
  t = \frac58.
\]
Define the base cutoff $z_{\mathrm{base}}:[0,1]^2\to[0,1]$ by
\[
 z_{\mathrm{base}}(x,y)=
 \begin{cases}
   1, & y\ge t\ \text{and }x<t,\\
   t, & y\ge t\ \text{and }x\ge t,\\
   0, & y<t\ \text{and }x\ge t,\\
   t, & y<t,\ x<t,\ x\le y,\\
   y, & y<t,\ x<t,\ x>y.
 \end{cases}
\]
Inside the square $[t_1,t]^2$, replace this by
\[
 z_0(x,y)=
 \begin{cases}
   t_1, & y<t_2,\ x\ge t_2,\\
   t_2, & y<t_2,\ x<t_2,\ x\le y,\\
   y,   & y<t_2,\ x<t_2,\ x>y,\\
   t,   & y\ge t_2,\ x<t_2,\\
   t_2, & y\ge t_2,\ x\ge t_2,
 \end{cases}
\]
and outside $[t_1,t]^2$ set $z_0=z_{\mathrm{base}}$.  Finally define
\[
  f_0(x,y,z):=\1\{z<z_0(x,y)\}.
\]
This is the same rational-parameter rule formalized by Flin et
al.~\cite{flin2026lean}.

For a finite size $n$, use midpoint representatives
\[
  x_i:=\frac{2i+1}{2n},
  \qquad i\in[n],
\]
and define the midpoint discretization
\[
  h^{(n)}_0(i,j,k):=f_0(x_i,x_j,x_k).
\]
All comparisons against $t_1,t_2,t$ can be implemented using integer
cross-products, so no floating-point arithmetic is needed.

\begin{proposition}[Strict improvement without local flips]
For $n = 281$, the midpoint coloring $h^{(281)}_0$ has
\[
  M\bigl(h^{(281)}_0\bigr)=1{,}503{,}704{,}985.
\]
Consequently,
\[
  p^*\le
  \frac{1{,}503{,}704{,}985}{6{,}234{,}839{,}521}
  =0.2411778169967753\ldots
  <\frac{94835}{393216}.
\]
\end{proposition}

\begin{proof}
The exact checker evaluates the midpoint rule and counts all directed edges
by the aggregation identity proved in Section~\ref{sec:verification}.  It
returns
\[
  M=1{,}503{,}704{,}985,
  \qquad
  |E|=281^4=6{,}234{,}839{,}521.
\]
The strict comparison is certified by
\[
  1{,}503{,}704{,}985\cdot393216
  -94835\cdot6{,}234{,}839{,}521
  =-146{,}592{,}275<0.
\]
Apply the finite De Bruijn certificate lemma.
\end{proof}

\section{Refined certificate at n = 497}

The strongest certificate reported here uses the same midpoint base coloring
at $n = 497$ and then toggles a finite set of vertices.

Let
\[
  F\subseteq[497]^3,
  \qquad |F|=4750,
\]
be the explicit list supplied in the machine-readable certificate
\texttt{n497\_certificate.json}.  Define
\[
  h_{497}(v):=h^{(497)}_0(v)\oplus\1\{v\in F\}.
\]
The full list is deliberately kept outside the manuscript; the theorem
depends only on that fixed finite file and the exact checker.

\paragraph{Search used for the reported refinement.}
For reproducibility, the artifact package also contains the deterministic
search that generated $F$.  Let $W=\{218,219,\ldots,279\}$.  Starting from
$h^{(497)}_0$, for every candidate $v\in W^3$ let $\Delta M(v)$ be the exact
change in the monochromatic-edge count produced by toggling $v$.  At each
step, the search chooses a vertex with minimum negative $\Delta M(v)$,
breaking ties lexicographically by $(a,b,c)$, applies that toggle, and
repeats.  It stops when no negative change remains.  All changes are computed
with integer counts.  This procedure terminates after exactly $4750$ flips;
the cumulative change is $-77{,}600$, and the best remaining change is $0$.
The resulting ordered flip list reproduces the supplied certificate exactly.
This search description is included for experimental reproducibility and is
not needed for the validity of the fixed finite certificate.

For context, before applying the flips, the midpoint coloring already has
\[
  M\bigl(h^{(497)}_0\bigr)=14{,}715{,}071{,}917,
\]
corresponding to
\[
  0.2411775249912064\ldots,
\]
which is itself below $p_0$.  The $4750$ flips reduce the monochromatic-edge
count by exactly $77{,}600$.

\begin{theorem}[Improved finite upper bound]
The coloring $h_{497}$ has exactly
\[
  M(h_{497})=14{,}714{,}994{,}317
\]
monochromatic directed edges.  Therefore
\[
  \boxed{
  p^*\le
  \frac{14{,}714{,}994{,}317}{61{,}013{,}446{,}081}
  =0.2411762531404098\ldots
  }
\]
and, in particular,
\[
  \frac{14{,}714{,}994{,}317}{61{,}013{,}446{,}081}
  <\frac{94835}{393216}.
\]
\end{theorem}

\begin{proof}
The exact checker evaluates the fixed coloring $h_{497}$ and returns
\[
\begin{aligned}
  |V_{497}|&=497^3=122{,}763{,}473,\\
  M(h_{497})&=14{,}714{,}994{,}317,\\
  |E_{497}|-M(h_{497})&=46{,}298{,}451{,}764,\\
  |E_{497}|&=497^4=61{,}013{,}446{,}081.
\end{aligned}
\]
Hence the monochromatic-edge fraction is the displayed rational number.
The exact cross-product comparison with the published formally verified
value is
\[
\begin{aligned}
&14{,}714{,}994{,}317\cdot393216\\
&\qquad
-94{,}835\cdot61{,}013{,}446{,}081
=-38{,}953{,}738{,}163<0.
\end{aligned}
\]
The finite De Bruijn certificate lemma now gives the claimed upper bound on
$p^*$.
\end{proof}

The numerical improvement over $p_0$ is
\[
  \frac{94835}{393216}
  -\frac{14{,}714{,}994{,}317}{61{,}013{,}446{,}081}
  =\frac{38{,}953{,}738{,}163}
  {23{,}991{,}463{,}214{,}186{,}496}
  \approx1.6236499548\times10^{-6}.
\]
It is small, but exact.

\section{Exact verification}
\label{sec:verification}

A literal traversal of all $n^4$ directed edges is unnecessary.  For each
ordered pair $(b,c)\in[n]^2$, define
\[
  L_1(b,c):=\sum_{a\in[n]} h(a,b,c),
  \qquad
  R_1(b,c):=\sum_{d\in[n]} h(b,c,d).
\]
Among the $n^2$ directed edges of the form
\[
  (a,b,c)\to(b,c,d)
\]
with fixed $(b,c)$, the number of monochromatic edges is
\[
  L_1(b,c)R_1(b,c)
  +\bigl(n-L_1(b,c)\bigr)\bigl(n-R_1(b,c)\bigr).
\]
Similarly, the cut-edge count is
\[
  L_1(b,c)\bigl(n-R_1(b,c)\bigr)
  +\bigl(n-L_1(b,c)\bigr)R_1(b,c).
\]
Summing over $(b,c)$ yields the exact total in $O(n^3)$ arithmetic
operations.

\begin{lemma}[Aggregation identity]
For every coloring $h:[n]^3\to\{0,1\}$,
\[
  M(h)=
  \sum_{b,c\in[n]}
  \left[
    L_1(b,c)R_1(b,c)
    +(n-L_1(b,c))(n-R_1(b,c))
  \right].
\]
\end{lemma}

\begin{proof}
Fix $(b,c)$.  A directed edge $(a,b,c)\to(b,c,d)$ is monochromatic with
color $1$ exactly when both endpoints have color $1$, giving
$L_1R_1$ choices.  It is monochromatic with color $0$ exactly when both
endpoints have color $0$, giving $(n-L_1)(n-R_1)$ choices.  These cases are
disjoint and exhaustive among monochromatic edges with the fixed middle pair
$(b,c)$.  Summing over all $n^2$ middle pairs counts every directed edge
exactly once.
\end{proof}

The supplied C++ checker uses unsigned 64-bit integers for the edge counts and
128-bit integers for rational cross-products.  At $n = 497$, all raw counts are
well below $2^{64}$.  The final run verifies both
\[
  M+\mathrm{cut}=n^4
\]
and the exact target count before returning success.

\begin{table}[t]
\centering
\caption{Exact values relevant to the upper-bound improvement.  The last
column is $M\cdot393216-94835\cdot n^4$; a negative value certifies a strict
improvement over $94835/393216$.}
\begin{tabular}{lrrrr}
\toprule
Certificate & $n$ & Monochromatic edges $M$ & Fraction & Cross product \\
\midrule
Midpoint base & 281 & 1,503,704,985 & 0.2411778169967753 & -146,592,275 \\
Midpoint base & 497 & 14,715,071,917 & 0.2411775249912064 & -8,440,176,563 \\
Refined & 497 & 14,714,994,317 & 0.2411762531404098 & -38,953,738,163 \\
\bottomrule
\end{tabular}
\end{table}

\section{Reproducibility and certificate boundary}

\paragraph{Artifact availability.}
The certificate files listed below are maintained in the public GitHub
repository
\href{https://github.com/yonezaemon1-hub/finite-de-bruijn-certificate-one-round-2-coloring}{\texttt{yonezaemon1-hub/finite-de-bruijn-certificate-one-round-2-coloring}}.
The archival identifiers for this release are Paper DOI
\href{https://doi.org/10.5281/zenodo.22670656}{10.5281/zenodo.22670656}
and Software DOI
\href{https://doi.org/10.5281/zenodo.22669456}{10.5281/zenodo.22669456}.

The certificate package contains:
\begin{enumerate}[leftmargin=1.8em]
  \item \texttt{n497\_certificate.json}, containing the base-rule metadata and
  the explicit list of $4750$ flipped vertices;
  \item \texttt{reproduce\_n497\_search.cpp}, the standalone deterministic
  local-search reproducer that starts from the midpoint base coloring and
  regenerates the ordered flip list without reading the certificate;
  \item \texttt{SEARCH\_REPRODUCTION.md}, documenting the exact one-flip
  objective change, tie-breaking rule, stopping rule, and update scheme;
  \item \texttt{check\_n497\_embedded\_certificate.cpp}, a standalone exact
  checker with the flip list embedded in source form;
  \item \texttt{check\_n281\_base.cpp}, a standalone checker for the simpler
  no-flip certificate;
  \item \texttt{reproduce\_all.sh} and \texttt{reproduce\_all.ps1}, which
  compile the reproducer and checkers, regenerate the $n=497$ certificate,
  compare it byte-for-byte with the canonical JSON, and rerun both exact
  checkers;
  \item logs recording the exact integer outputs of the search and checkers.
\end{enumerate}

For the fixed files used in this draft, the SHA-256 values are
\begin{quote}\small\ttfamily
n497\_certificate.json\\
41e62c73421e9ca33b097ecedda6d31b\\
398ed8189d74444258c8015e197f193c\\[0.5ex]
reproduce\_n497\_search.cpp\\
4bb0bfe5467deb76a5e2861b0c13c2f\\
8a95de6e612ecd7d11ce5d2bbbfa7a3db\\[0.5ex]
check\_n497\_embedded\_certificate.cpp\\
06280be157851220138e8fec38d2852ac\\
da1b59543a0838040d57a3e171cf582
\end{quote}
Each two-line hexadecimal value is to be concatenated without whitespace.

The theorem itself still requires only reconstructing the fixed finite
coloring, evaluating the aggregation identity, and checking one integer
inequality.  The additional search reproducer closes the stronger experimental
reproducibility loop: base coloring $\to$ deterministic local search $\to$
regenerated certificate $\to$ exact checker.  The reproduced JSON is
byte-identical to the canonical certificate.  No stochastic optimization,
floating-point tolerance, or solver output is involved in either the search
reproduction or the proof-certificate verdict.

\section{Discussion and limitations}

The result improves the published exact value of the three-parameter
construction of Flin et al., but the improvement is numerically modest.  It narrows the interval between the published lower
bound $0.23879$ and the upper bound by about $0.068\%$ relative to the
previous interval width.

More importantly, the result should not be interpreted as evidence that the
finite coloring is close to optimal.  We have not computed
$p_{\normal}(497)$ exactly, nor do we provide a lower bound on the maximum cut
of $\DB_{\normal}(497)$.  Further finite local search, stronger combinatorial
optimization, or a new structured continuous rule may lower the upper bound
again.

The $n = 281$ proposition also highlights a discretization effect: the finite
midpoint evaluation of the existing continuous rule can lie slightly below
its continuous exact value.  This is fully legitimate in the De Bruijn
framework because each finite coloring independently induces a valid
one-round randomized algorithm through the discretization lemma.  The
$n = 497$ local refinement then exploits the additional freedom of arbitrary
finite colorings beyond direct samples of the continuous surface.

\section{Conclusion}

We have exhibited a finite, exactly checkable De Bruijn coloring giving
\[
  p^*\le0.2411762531404098\ldots,
\]
strictly improving the formally verified value
$94835/393216=0.2411778767903646\ldots$ associated with the previously
published three-parameter rule.  The theorem follows from an explicit
$497$-symbol coloring, exact integer edge counting, and the normal De Bruijn
upper-bound lemma of Flin et al.  A simpler $281$-symbol midpoint
discretization already suffices for a strict improvement, while the refined
$497$-symbol certificate yields the stronger numerical bound reported here.

\begin{thebibliography}{9}

\bibitem{flin2026podc}
M.~Flin, A.~Raevskaya, R.~Stimpert, J.~Suomela, and Q.~Yang.
\newblock Brief Announcement: 2-Coloring Cycles in One Round.
\newblock In \emph{Proceedings of the 2026 ACM Symposium on Principles of
Distributed Computing (PODC 2026)}, pages 44--47, 2026.
\newblock DOI: \href{https://doi.org/10.1145/3796701.3815937}{10.1145/3796701.3815937}.

\bibitem{flin2026arxiv}
M.~Flin, A.~Raevskaya, R.~Stimpert, J.~Suomela, and Q.~Yang.
\newblock 2-Coloring Cycles in One Round.
\newblock arXiv:2603.04235v2, March 2026.
\newblock \href{https://arxiv.org/abs/2603.04235}{arXiv:2603.04235}.

\bibitem{flin2026lean}
M.~Flin, A.~Raevskaya, R.~Stimpert, J.~Suomela, and Q.~Yang.
\newblock 2-Coloring Cycles in One Round: Formalization in Lean 4.
\newblock GitHub repository \texttt{suomela/2-coloring-1-round}, 2026.
\newblock Pinned commit: \texttt{aa04b54db9a48d7a5084e263c147607af7744498}.
\newblock Parameter definitions:
\texttt{Distributed2Coloring/UpperBound/Recursive3Param.lean}.
\newblock Exact rational value:
\texttt{Distributed2Coloring/UpperBound/Recursive3Param/Final.lean}.

\end{thebibliography}

\end{document}
